\section{P04 - Fiche cours - Dictionnaires}

\subsection{À savoir}

\begin{itemize}
\tightlist
\item
  types construits se travaille dans le contexte ``tuples, listes et
  dictionnaires'' avec des données vérifiables.
\item
  La fiche distingue vocabulaire, méthode, exemple corrigé et contrôle
  pour dictionnaires.
\item
  Les capacités P-DATA-CONSTR-03A, P-DATA-CONSTR-03B, P-DATA-CONSTR-03C
  sont rappelées ici sans être déclarées couvertes.
\item
  L'élève doit pouvoir refaire un exemple de types construits avec une
  valeur, une table ou un code différent.
\end{itemize}

\subsection{Méthodes}

\begin{enumerate}
\def\labelenumi{\arabic{enumi}.}
\tightlist
\item
  Capacités explicitement travaillées dans les méthodes et exercices :
  P-DATA-CONSTR-03A, P-DATA-CONSTR-03B, P-DATA-CONSTR-03C.
\item
  P-DATA-CONSTR-03A : choisir une structure selon l'accès attendu.
\item
  Identifier les données d'entrée de dictionnaires puis écrire le
  résultat attendu avant de conclure.
\item
  Contrôler dictionnaires par un cas limite explicite et une
  vérification courte.
\item
  Relier la réponse à un support de séance P04 sans confondre fiche de
  révision et preuve de couverture.
\end{enumerate}

\subsection{Exemples corrigés}

\subsubsection{Exemple corrigé 1 - Exemple principal}

\texttt{notes{[}1{]}\ =\ 14} modifie la deuxième valeur de
\texttt{{[}8,12,10{]}}. \#\#\# Exemple corrigé 2 - Contrôle ou
contre-exemple \lstinline|eleve["score"]| lit une valeur par clé dans
un dictionnaire.

\subsection{Erreurs fréquentes}

\begin{itemize}
\tightlist
\item
  Confondre le vocabulaire de dictionnaires avec une simple récitation :
  corriger par un exemple calculé ou exécuté.
\item
  Oublier une hypothèse de tuples, listes et dictionnaires : corriger en
  l'écrivant avant la méthode.
\item
  Conclure sans contrôle sur types construits : corriger par un cas
  limite ou une vérification inverse.
\end{itemize}

\subsection{Cas limites}

\begin{itemize}
\tightlist
\item
  Cas de départ vide ou nul pour dictionnaires, à traiter selon la
  convention du chapitre P04.
\item
  Donnée invalide dans tuples, listes et dictionnaires, par exemple
  symbole interdit, clé absente ou requête trop large selon la fiche.
\item
  Cas frontière de types construits où une seule valeur change la
  méthode ou le résultat attendu.
\end{itemize}

\subsection{Mini-exercices}

\subsubsection{Mini-exercice 1}

P-DATA-CONSTR-03A : appliquer la méthode de dictionnaires à un exemple
court choisi dans le chapitre P04. \#\#\# Mini-exercice 2 Repérer
l'erreur dans une réponse qui oublie une hypothèse de tuples, listes et
dictionnaires. \#\#\# Mini-exercice 3 Proposer un cas limite pertinent
pour types construits et expliquer le résultat attendu. \#\#\#
Mini-exercice 4 Écrire une phrase de contrôle qui vérifie la conclusion
obtenue pour dictionnaires.

\subsection{Réponses rapides}

\begin{enumerate}
\def\labelenumi{\arabic{enumi}.}
\tightlist
\item
  La méthode attendue pour dictionnaires commence par les données puis
  applique l'opération du chapitre P04.
\item
  L'erreur vient de l'hypothèse manquante ; elle se corrige en testant
  le cas mentionné dans tuples, listes et dictionnaires.
\item
  Le cas limite doit donner un résultat explicite, par exemple 0, vide,
  absent ou hors plage selon types construits.
\item
  Le contrôle compare le résultat avec la définition ou avec une
  opération inverse de dictionnaires.
\end{enumerate}

\subsection{À retenir}

\begin{itemize}
\tightlist
\item
  P04 : dictionnaires se révise avec une définition, une méthode et un
  exemple corrigé.
\item
  Les capacités P-DATA-CONSTR-03A, P-DATA-CONSTR-03B, P-DATA-CONSTR-03C
  restent en travail tant que TD, TP, évaluation, barème et revues
  humaines manquent.
\item
  Un exemple de types construits doit changer autre chose qu'une simple
  valeur pour tester la compréhension.
\item
  Pour P04, le tableau de liens distingue les supports existants et les
  supports inscrits au registre.
\item
  La fiche P04 sur dictionnaires reste needs\_review et ne déclenche ni
  publication ni couverture.
\end{itemize}

\subsection{Lien avec la progression}

\begin{longtable}[]{@{}
  >{\raggedright\arraybackslash}p{(\columnwidth - 6\tabcolsep) * \real{0.2500}}
  >{\raggedright\arraybackslash}p{(\columnwidth - 6\tabcolsep) * \real{0.2500}}
  >{\raggedright\arraybackslash}p{(\columnwidth - 6\tabcolsep) * \real{0.2500}}
  >{\raggedright\arraybackslash}p{(\columnwidth - 6\tabcolsep) * \real{0.2500}}@{}}
\toprule\noalign{}
\begin{minipage}[b]{\linewidth}\raggedright
Élément
\end{minipage} & \begin{minipage}[b]{\linewidth}\raggedright
Fichier
\end{minipage} & \begin{minipage}[b]{\linewidth}\raggedright
Statut
\end{minipage} & \begin{minipage}[b]{\linewidth}\raggedright
Remarque
\end{minipage} \\
\midrule\noalign{}
\endhead
\bottomrule\noalign{}
\endlastfoot
Séance & P04-S1 & prête & séance présente dans la progression \\
TD &
03\_progressions/supports/premiere/P04/P04\_td\_types\_construits.md &
existant & support associé existant dans 03\_progressions/supports \\
TP &
03\_progressions/supports/premiere/P04/P04\_tp\_types\_construits.md &
existant & support associé existant dans 03\_progressions/supports \\
Évaluation &
03\_progressions/supports/premiere/P04/P04\_evaluation\_types\_construits.md
& existant & support associé existant dans 03\_progressions/supports \\
\end{longtable}

\subsection{Auto-évaluation}

\begin{itemize}
\tightlist
\item
  Je peux expliquer dictionnaires avec un exemple différent de ceux de
  la fiche P04.
\item
  Je peux citer au moins une capacité parmi P-DATA-CONSTR-03A,
  P-DATA-CONSTR-03B, P-DATA-CONSTR-03C et dire où elle est travaillée
  dans la fiche.
\item
  Je peux dire quel support lié à P04 existe déjà ou reste inscrit au
  registre.
\item
  Je peux identifier un cas limite de types construits sans transformer
  la fiche en corrigé complet.
\end{itemize}
